\section{Why Axioms?}
\label{sec:why-axioms}

% Cross-references to chapters 1-3 (sec:and-under-noise, sec:composition-theorem)
% will resolve when those chapters are drafted in Plans 2-4. Cross-references
% to sibling sections sec:axiom-1, sec:axiom-2, sec:kronecker, and
% sec:bernoulli-bool-formal will resolve as Tasks 3-6 of this plan land.
% Build warnings ("Reference X undefined") for these labels are expected
% during ch 4 drafting and will be resolved by Plan 4's Part I integration
% check.

The three chapters behind us built something real. We derived the error
rates for AND and OR gates operating on noisy bits
(\Cref{sec:and-under-noise}), showed that composition of
channels produces predictable compound rates, and in \Cref{sec:composition-theorem}
proved the composition theorem: when $k$ independent noisy channels
are chained in series, their true-positive rates compose multiplicatively.
We then applied that theorem immediately to Miller--Rabin primality testing,
watching a per-round false-positive rate of at most $1/4$ collapse to
$(1/4)^k$ after $k$ rounds. The derivations were concrete and the results
were correct.

But correct derivations can hide assumptions. Chapters 1--3 made
\emph{two} independence assumptions so naturally that we did not pause
to name them. This chapter names them, states them precisely, and then
shows exactly what mathematical structure they license. That structure
is a precise structural theorem showing what those assumptions cost and
what they buy; it is the formal payoff for the care we are about to take.

\subsection*{The first hidden assumption: element-wise independence}

Recall the composition theorem from \Cref{sec:composition-theorem}.
When we wrote the true-negative rate for $k$ channels chained in series
as $\tnrate_{\mathrm{total}} = 1 - \prod_{i=1}^k (1 - \tnrate_i)$, we
were treating each channel's contribution to the end-to-end error as
independent of every other channel's contribution. One query enters the
chain; it passes through channel 1, whose output feeds channel 2, and so
on; the overall error rate is a product of independent per-stage terms.
No correlation across stages was allowed, and the product formula relied
on that fact throughout.

A different scenario where Axiom 1 matters: two classifiers trained on
overlapping data from the same source. Their errors on a new input $x$
will tend to correlate: when one misclassifies $x$, the other is more
likely to as well. The product formula breaks. The effective error rate
is worse than the product formula predicts, because the errors reinforce
rather than cancel.

We will call the assumption that justifies the product formula
\emph{Axiom 1}. Roughly, Axiom 1 asserts independence of error events
across distinct input elements; the precise statement and its scope
appear in \Cref{sec:axiom-1}. The practical effect is that joint error
distributions factor across distinct queries, which is what made the
composition theorem's product form possible. When Axiom 1 holds, every
formula in chapters 1--3 that relies on independence is automatically
justified. When Axiom 1 fails, those formulas become bounds at best and
misleading at worst. Chapter 3's Miller--Rabin analysis, for instance,
requires Axiom 1: the $k$ witness tests must be independent for the
per-round bound to compound as $(1/4)^k$. If the witnesses were
correlated, the product would underestimate the true composite error
rate.

\subsection*{The second hidden assumption: stable block error rates}

The derivation of AND under noise in \Cref{sec:and-under-noise} used
specific values $\fprate$ and $\fnrate$, the false-positive and
false-negative rates, treating them as fixed properties of each input.
We implicitly assumed that querying an input repeatedly would yield the
same error behavior each time: that $\fprate$ does not drift from query
to query, that the channel is stationary across uses.

This is not the only conceivable behavior. A Markov-correlated channel,
for instance, lets the error rate on query $t+1$ depend on what
happened at query $t$. If the channel just made a false positive, it
might be more likely to make another. The rates are not stable; they are
history-dependent. We did not model this, and for good reason: the
stationary case is the tractable and practically important one. But
we should say we are assuming it.

We have been assuming Axiom 2 since \Cref{sec:and-under-noise}; let us
make it explicit. Axiom 2 says: the error rates within each partition
block of the input domain are the same for every query, unconditionally.
There is one channel matrix $Q$ per input-block-pair, and it does not
depend on which query in a sequence we are currently processing.
Under Axiom 2, the channel is genuinely memoryless within each block:
its statistical behavior is the same on the hundredth query as on the
first.

This is a meaningful restriction, not a triviality. It rules out
fatigue, warm-up effects, adversarial adaptation, and history
dependence. It also makes the math work cleanly. We formalize Axiom 2
in \Cref{sec:axiom-2}.

\subsection*{What the axioms buy}

Once both axioms are in place, something striking follows. The joint
confusion matrix describing $m$ independent queries on the channel is
not a free matrix with exponentially many parameters. Under Axioms 1
and 2, it \emph{must} factor into a product of $m$ small per-query
matrices, one factor per query. The exponential-to-linear collapse in
parameter count, made precise in \Cref{sec:kronecker}, is the formal
payoff of this chapter.

This is not merely an algebraic convenience. It is the theorem that
makes the Bernoulli framework useful in practice. Without it, a
fifty-query analysis would require specifying an astronomically large
matrix. With it, you need fifty pairs of numbers. The saving is
exponential.

We prove this Kronecker factorization theorem, derive the parameter
count corollary, and work through explicit examples in
\Cref{sec:kronecker}. The theorem is the chapter's climax and the
formal vindication of every computation in chapters 1--3.

\subsection*{The formal definition}

Once the axioms and the Kronecker theorem are in hand, we can at last
write down the formal definition of the type $\Bernoulli{\BoolType}$
without fear of circularity. Chapter 1 named the noisy bit informally.
Chapter 2 gave it a channel matrix. Chapter 3 showed how to compose
these objects. Chapter 4 closes by showing that all of those things
are instances of a single precise definition: a triple $(X, Y, Q)$
where $X$ and $Y$ are the input and output domains $\{\top, \bot\}$,
$Q$ is a $2 \times 2$ stochastic matrix, and the whole structure
satisfies Axioms 1 and 2. We work this out in \Cref{sec:bernoulli-bool-formal},
including the biased-coin example from \S 1.3 that was deliberately
left without a full PMF until we had the formal machinery to handle it.

The trajectory of this chapter, then, is: axioms first
(\Cref{sec:axiom-1,sec:axiom-2}), then the theorem they make possible
(\Cref{sec:kronecker}), then the definition the theorem earns us
(\Cref{sec:bernoulli-bool-formal}).

By the end, the reader who worked through chapters 1--3 will recognize
every example, every formula, and every derivation as an instance of
one coherent definition. That is the point.
